% === §6 Metodología ===
% Redactado por el Redactor — Fase 5
% Estructura: cadena de valor por objetivo específico (F1→F2→F3→F4)
% Mapeo estricto: F1↔OE1↔SP1↔E1; F2↔OE2↔SP2↔E2+E3+E4;
%                F3↔OE3↔SP3↔E5 (TRL 6); F4↔OE3↔SP3↔E6 (proyección TRL 7)

% --- Párrafo introductorio: la metodología como cadena de valor ---

La metodología del proyecto se estructura como una \textbf{cadena de valor} cuyos
eslabones son los tres objetivos específicos (OE1$\to$OE2$\to$OE3) desarrollados
secuencialmente en cuatro fases encadenadas (F1$\to$F2$\to$F3$\to$F4), conforme al
alcance definido en la sección~\ref{sec:alcance}. Cada fase consume los
entregables de la anterior y produce insumos verificables para la siguiente: la
fase de fundamentación y caracterización del aprendiz (F1, OE1) entrega el modelo
computacional del aprendiz que habilita el diseño del ecosistema agéntico
(F2, OE2); éste, una vez integrado y validado en entorno simulado, alimenta la
validación empírica en aula (F3, OE3), cuyas evidencias de impacto sustentan la
transferencia tecnológica y la proyección TRL~7 (F4, OE3). El punto de partida del
equipo se sitúa en TRL~3--4 (prototipos de modelos individuales validados en
laboratorio \cite{escamilla2022artificial}); el proyecto avanza dos niveles de
madurez hasta alcanzar TRL~6 al cierre de F3 y proyecta TRL~7 en F4. Cada fase
define actividades, personal responsable, sustento teórico-metodológico (enfoques
E1--E6 de \S5.3), recursos tecnológicos y experiencia previa del equipo.

\medskip

% ===== FASE 1: Fundamentación y caracterización del aprendiz (OE1, SP1, E1) =====
\noindent\textbf{Fase 1 — Fundamentación y caracterización del aprendiz
(OE1, SP1, Enfoque 1).}
\emph{Responsable principal:} Investigador principal (IP, Andrés Álvarez) con
coinvestigador especialista en psiquiatría educativa / matemáticas (David
Angulo); apoyo de un estudiante de doctorado y un estudiante de pregrado del
semillero. \emph{Duración estimada:} meses 1--5 (5 meses).

\emph{Actividades:} (i) Diseño y validación psicométrica del instrumento de
diagnóstico de aprendizaje profundo —motivación intrínseca, argumentación
técnica y resolución de problemas auténticos— mediante juicio de expertos
(IVC $\geq 0.80$) y un estudio piloto con estudiantes voluntarios (alfa de Cronbach
$\geq 0.80$). (ii) Recolección y curaduría de un \emph{dataset} multimodal de
trazas (textos argumentativos, soluciones de código, registros de simulación,
interacciones con IA) en asignaturas reales de los programas objetivo, bajo
protocolos de consentimiento informado aprobados por el Comité de Ética de la
Facultad de Ingeniería y Arquitectura. (iii) Desarrollo y entrenamiento del modelo
computacional del aprendiz (Enfoque 1) mediante transformers en español (BETO,
RoBERTa-BNE) y clasificadores supervisados, con validación cruzada estratificada
$k$-fold ($k=5$) reportando F1-score macro y AUC-ROC.

\emph{Sustento teórico-metodológico (Enfoque 1):} modelado computacional
multidimensional del aprendiz que combina NLP basado en transformers y
aprendizaje automático supervisado para inferir el nivel de aprendizaje profundo
a partir de trazas multimodales, operacionalizando los constructos de Marton \&
Säljö y el modelo SOLO de Biggs en señales computacionalmente trazables
\cite{hooshyar2024systematic}.

\emph{Recursos tecnológicos:} estaciones Apple Silicon del LabIA para inferencia
\emph{on-device} de modelos preentrenados; repositorio de datos con gobernanza
ética; \emph{frameworks} HuggingFace Transformers, scikit-learn, XGBoost.

\emph{Experiencia previa:} el equipo cuenta con trabajo previo en analítica del
aprendizaje con \emph{deep learning} y NLP aplicado a educación en ingeniería
\cite{escamilla2022artificial}, lo que sitúa la línea base de los componentes
individuales en TRL~3--4.

\emph{Hito de fase H1 (mes 5):} instrumento validado, \emph{dataset} curado y
modelo del aprendiz entregados; nivel de madurez TRL~4 confirmado para los
componentes de diagnóstico.

\medskip

% ===== FASE 2: Diseño y desarrollo del ecosistema agéntico (OE2, SP2, E2+E3+E4) =====
\noindent\textbf{Fase 2 — Diseño y desarrollo del ecosistema agéntico
(OE2, SP2, Enfoques 2, 3 y 4).}
\emph{Responsable principal:} IP con coinvestigador en arquitecturas de software
(César Castellanos); \emph{team lead} de desarrollo e ingenieros del semillero
(estudiantes de maestría en Ing.\ Automática y pregrado en Cs.\ de la Computación /
Ing.\ de Software). \emph{Duración estimada:} meses 4--12 (9 meses, con solape
inicial con F1 para arrancar el diseño arquitectónico mientras concluye la
caracterización del aprendiz).

\emph{Actividades:} (i) Diseño de la \emph{ontología de roles agénticos}
—tutor, evaluador, generador de problemas, simulador de sistemas
eléctricos/electrónicos y asistente de programación— con fundamento pedagógico
explícito (Enfoque 2). (ii) Implementación del \emph{bus de orquestación} basado
en el Model Context Protocol (MCP), con memoria compartida y contexto
pedagógico que consume el modelo del aprendiz producido en F1. (iii) Desarrollo
del módulo de \emph{scaffolding} adaptativo (Enfoque 3), que define la función
$s(\mathbf{x}) \to \{h, d, c\}$ que ajusta los parámetros heurística,
direccionalidad y complejidad según el perfil del aprendiz. (iv) Implementación
del \emph{generador de problemas auténticos} y del \emph{evaluador formativo}
(Enfoque 4) mediante tuberías RAG sobre una base documental de ingeniería y NLP
para \emph{automated short answer scoring}. (v) Integración del ecosistema y
pruebas en entorno simulado con estudiantes voluntarios para verificar
interoperabilidad, latencia ($p_{95}$), tasa de alucinación factual y acuerdo
interevaluador ($\kappa \geq 0.70$).

\emph{Sustento teórico-metodológico (Enfoques 2--4):} arquitectura multiagente
heterogénea con orquestación MCP (E2) que resuelve la fragmentación de
herramientas; scaffolding adaptativo computacional basado en la ZDP de
Vygotsky (E3) que gradúa la asistencia evitando sobreasistencia y subasistencia;
generación de problemas auténticos con RAG y evaluación formativa automatizada
(E4) que cierra el ciclo problema$\to$respuesta$\to$retroalimentación.

\emph{Recursos tecnológicos:} LLMs preentrenados (locales: Llama~3, Mistral;
vía API: GPT-4o, Claude~3.5); SDKs MCP en Python/TypeScript; \emph{frameworks}
LangGraph, LangChain, LlamaIndex; simuladores SPICE, intérpretes Python y
entornos de compilación integrados vía MCP; VPS Hostinger para orquestación en
nube.

\emph{Experiencia previa:} el equipo ha evaluado \emph{frameworks} agénticos y
protocolos de orquestación en proyectos del semillero, y cuenta con
familiaridad con \emph{spec-driven development} y \emph{vibe coding} en
infraestructura Apple Silicon.

\emph{Hito de fase H2 (mes 12):} prototipo integrado del ecosistema agéntico
validado en entorno simulado; nivel de madurez TRL~5 alcanzado.

\medskip

% ===== FASE 3: Validación empírica en aula (OE3, SP3, E5) — TRL 6 =====
\noindent\textbf{Fase 3 — Validación empírica en aula del impacto causal
(OE3, SP3, Enfoque 5).}
\emph{Responsable principal:} IP con coinvestigador especialista en psicometría
y educación (David Angulo); un estudiante de doctorado y dos de pregrado
encargados del trabajo de campo y la analítica. \emph{Duración estimada:} meses
11--16 (6 meses, con solape inicial para el diseño del protocolo mientras
concluye la integración del ecosistema).

\emph{Actividades:} (i) Diseño del protocolo cuasiexperimental pre--post con
grupo de control no equivalente (Enfoque 5), selección de asignaturas
—Circuitos Eléctricos, Programación Orientada a Objetos, Sistemas Embebidos—,
definición de grupos experimental y control ($n \geq 30$ c/u), y aprobación del
Comité de Ética. (ii) Ejecución de la intervención a lo largo de un semestre
académico: medición basal (T0), uso del ecosistema agéntico por el grupo
experimental mientras el grupo de control recibe instrucción tradicional con
herramientas de IA convencionales, y medición final (T2). (iii) Análisis
estadístico: ANOVA de medidas repetidas (factor grupo $\times$ tiempo), tamaño
del efecto ($d$ de Cohen, $\eta^2$ parcial) y ANCOVA para controlar diferencias
basales, complementado con analítica del aprendizaje (relación
dosis--respuesta). (iv) Validación cualitativa: entrevistas semiestructuradas,
análisis temático de la argumentación técnica y grupos focales con docentes.

\emph{Sustento teórico-metodológico (Enfoque 5):} el diseño cuasiexperimental
permite aislar el efecto atribuible al ecosistema controlando estadísticamente
las diferencias basales y los factores de confusión (maduración, Hawthorne,
regresión a la media), produciendo estimaciones del tamaño del efecto comparables
con la literatura internacional sobre intervenciones educativas efectivas
\cite{freeman2014active}.

\emph{Recursos tecnológicos:} instancia del ecosistema desplegada para el aula;
plataforma de \emph{learning analytics} para captura de trazas; instrumentos
psicométricos de F1; software estadístico (R, Python \texttt{statsmodels}).

\emph{Hito de fase H3 (mes 16):} evidencia cuantitativa y cualitativa del
impacto causal del ecosistema en el aprendizaje profundo; verificación del
nivel de madurez TRL~6 mediante lista de cotejo estandarizada (demostración en
entorno relevante —aula universitaria real—).

\medskip

% ===== FASE 4: Transferencia, despliegue y proyección TRL 7 (OE3, SP3, E6) =====
\noindent\textbf{Fase 4 — Transferencia tecnológica, despliegue y proyección
TRL~7 (OE3, SP3, Enfoque 6).}
\emph{Responsable principal:} IP con el \emph{team lead} de desarrollo;
coinvestigadores docentes apoyan la redacción y diseminación; un profesional de
transferencia tecnológica y un estudiante de posgrado gestionan la adopción
externa. \emph{Duración estimada:} meses 15--20 (6 meses, con solape inicial
para iniciar el despliegue mientras concluyen los análisis de F3; los meses
19--20 corresponden al periodo adicional de 2 meses para productos estipulado en
la convocatoria).

\emph{Actividades:} (i) Despliegue del ecosistema en infraestructura híbrida
—estaciones Apple Silicon locales para inferencia \emph{on-device} y VPS en nube
para orquestación y analítica (Enfoque 6). (ii) Documentación de
interoperabilidad (API, MCP, estándares LTI/Caliper) y de procedimientos de
instalación y operación; obtención del registro de software ante la Dirección
Nacional de Derechos de Autor de Colombia; publicación del código bajo licencia
abierta permisiva en GitHub. (iii) Ejecución del plan de adopción con al menos
un actor del sector productivo (empresa de software o centro de formación técnica
del eje cafetero) que demuestre la operación del ecosistema en un entorno
relevante. (iv) Redacción y sometimiento de artículos en revistas indexadas
Q1/Q2, ponencias y productos de apropiación social del conocimiento.

\emph{Sustento teórico-metodológico (Enfoque 6):} el marco TRL adaptado a
\emph{edtech} proporciona criterios operacionales verificables (TRL~6:
demostración en entorno relevante; TRL~7: demostración operativa en entorno del
sector productivo) que guían la maduración y generan evidencia documentada para
los potenciales adoptantes, reduciendo la incertidumbre de transferencia.

\emph{Recursos tecnológicos:} infraestructura híbrida de despliegue; repositorios
GitHub; canal de adopción con el actor externo; plantillas de documentación TRL.

\emph{Hito de fase H4 (mes 20):} ecosistema desplegado y documentado, registro de
software obtenido, demostración operativa con actor externo y productos
académicos entregados; proyección TRL~7 documentada.

\medskip

% ===== Descripción del diagrama metodológico para el Diseñador =====
% NOTA PARA EL DISEÑADOR (agente TikZ): a continuación se describe el contenido
% conceptual. Implementar en proposal/sections/diag_metodologico.tex.

\noindent\textbf{Propuesta de contenido para el diagrama esquemático metodológico.}
El diagrama debe representar la cadena de valor del proyecto como un flujo
horizontal de cuatro fases consecutivas (F1$\to$F2$\to$F3$\to$F4), donde cada
fase es un bloque rectangular etiquetado con: el objetivo específico asociado
(OE1, OE2, OE3), el subproblema resuelto (SP1, SP2, SP3), los enfoques
sustentantes (E1; E2+E3+E4; E5; E6) y el rol responsable principal.
Sobre el flujo horizontal debe dibujarse una \textbf{cinta de TRL} paralela
que indique claramente la progresión de madurez: TRL~4 (punto de partida, al
inicio de F1), TRL~5 (cierre de F2), TRL~6 (cierre de F3, marcado como meta
intermedia destacada) y TRL~7 (cierre de F4, marcado como proyección). Las
flechas entre fases deben llevar etiquetados los entregables que transfieren de
un eslabón al siguiente —p.\,ej., ``modelo del aprendiz'' (F1$\to$F2),
``ecosistema integrado TRL~5'' (F2$\to$F3), ``evidencia causal de impacto''
(F3$\to$F4)—. En un panel inferior o lateral se deben destacar las
\textbf{novedades metodológicas y algorítmicas} clave: el modelo computacional
del aprendiz basado en transformers en español, la arquitectura multiagente con
orquestación MCP, el scaffolding adaptativo computacional y la tubería RAG de
generación de problemas auténticos con evaluación formativa automatizada. En un
panel opuesto (superior o lateral contrario) se deben representar los
\textbf{beneficiarios finales}: estudiantes de pregrado y posgrado de
ingeniería, docentes de la UNAL, instituciones aliadas vía SIUN y actores del
sector productivo colombiano (empresas de software y centros de formación
técnica). El diagrama debe usar una paleta cromática institucional (azul UNAL,
gris LabIA, verde GCPDS) y ser legible en escala de grises. El Diseñador deberá
asegurar que la correspondencia fase$\leftrightarrow$objetivo$\leftrightarrow$
subproblema$\leftrightarrow$enfoque sea trazable visualmente.